%=======================================================================
%  VernaSolver — IEEE Journal / Transactions style RESEARCH PAPER
%  How to use in Overleaf:
%    1. Create a new blank project.
%    2. Paste this whole file into main.tex.
%    3. Menu (top-left) → Settings → Compiler → pdfLaTeX.
%    4. Click "Recompile". IEEEtran is built into Overleaf.
%  This uses the [journal] mode of IEEEtran — the standard IEEE
%  research-paper (Transactions) format, not the conference format.
%  Layout, spacing and alignment are handled by the class — do not
%  change the page geometry.
%=======================================================================
\documentclass[journal]{IEEEtran}

% --- Packages ---------------------------------------------------------
\usepackage{cite}
\usepackage{amsmath,amssymb,amsfonts}
\usepackage{algorithmic}
\usepackage{graphicx}
\usepackage{textcomp}
\usepackage{xcolor}
\usepackage{url}
\usepackage{booktabs}
\usepackage{tikz}
\usetikzlibrary{shapes.geometric, arrows.meta, positioning, fit, backgrounds}

\def\BibTeX{{\rm B\kern-.05em{\sc i\kern-.025em b}\kern-.08em
    T\kern-.1667em\lower.7ex\hbox{E}\kern-.125emX}}

\begin{document}

% =====================================================================
%  TITLE
% =====================================================================
\title{A Textbook-Grounded Question Answering System with Page-Level
Citations and Multilingual Support for Students}

% ---------------------------------------------------------------------
%  AUTHORS  (IEEE journal/Transactions style — flat author line)
% ---------------------------------------------------------------------
\author{Pritam~Bairagi,~Shwetaki~Killedar,~and~Preksha~Pokharna%
\thanks{The authors are with the Department of Computer Science and
Engineering, MIT Art, Design and Technology University, Pune, India.}%
\thanks{Enrollment numbers: P. Bairagi (ADT24SOCB0846),
S. Killedar (ADT24SOCB1158), P. Pokharna (ADT24SOCB0839).}}

% Running headers for the journal style
\markboth{VernaSolver: Textbook-Grounded Question Answering with
Page-Level Citations}{Bairagi \MakeLowercase{\textit{et al.}}: A
Textbook-Grounded Question Answering System for Students}

\maketitle

% =====================================================================
%  ABSTRACT
% =====================================================================
\begin{abstract}
Students often use AI chat tools to study, but these tools answer from
general training data and not from the book a student is actually
following. This can lead to answers that do not match the course book,
page numbers that are not real, and weak support for languages like
Hindi and Marathi. This paper describes \textit{VernaSolver}, a system
that takes a student's own textbook and answers questions only from it,
giving the page number for each answer so the student can check it. The
book is uploaded as a PDF, split into small parts, and stored in a
searchable form. For each question, the system picks the most useful
parts of the book, orders them with a second model, and asks an AI model
to write the answer using only those parts. When the book does not
contain the answer, the system says so rather than making one up. The
tool can search one book or a whole subject, reply in English, Hindi, or
Marathi, and build short quizzes and flashcards from the book. We present
the design, the steps that keep answers tied to the source, and the
results from a working version tested on real course books.
\end{abstract}

\begin{IEEEkeywords}
Question Answering, Retrieval-Augmented Generation, Vector Search,
Textbook Study Tool, Page Citation, Multilingual, Education.
\end{IEEEkeywords}

\IEEEpeerreviewmaketitle

% =====================================================================
%  I. INTRODUCTION
% =====================================================================
\section{Introduction}
AI chat tools have become a common way for students to get quick help.
They can explain ideas and answer questions on many topics. The trouble
starts when a student has to learn from one set book for a course or an
exam. A normal chat tool does not read that book. It answers from the
large data it was trained on, so its reply may not match the book, the
page it mentions may not exist, and it may go outside the syllabus. When
an exam is based on a fixed book, an answer the student cannot check is
not of much use.

Language adds another gap. Many students in India are more comfortable in
Hindi or Marathi, but most chat tools are tuned for English. So a student
may get an answer that is both hard to verify and in the wrong language.

We built \textbf{VernaSolver} to deal with both problems. The rule it
follows is plain: answer only from the student's book, and always show
where the answer came from. To do this it uses Retrieval-Augmented
Generation (RAG)~\cite{lewis2020rag}, where the system first looks up the
right parts of the book and then asks the AI model to answer from those
parts only, with the page number shown for each one.

This paper makes the following points:
\begin{itemize}
\item It describes a study tool that answers only from a chosen book and
gives a page number that the student can check.
\item It explains a search step that finds parts by meaning and then
re-orders them with a second model for better accuracy, and rewrites
follow-up questions so longer chats work.
\item It shows a low-cost setup where the search runs on a normal
computer and the AI model is used only to write the final answer, with a
backup model if the main one is down.
\item It adds useful study features: subject-wide search, replies in
three languages, and automatic quizzes and flashcards, through both a
command-line tool and a web app.
\end{itemize}

% =====================================================================
%  II. RELATED WORK
% =====================================================================
\section{Related Work}
\textbf{Grounding answers in a source.} Lewis \textit{et
al.}~\cite{lewis2020rag} proposed RAG, where a generation model is paired
with a search step over an outside source so it does not rely only on
memory. This has been shown to cut down made-up answers and to make it
easier to point back to the source.

\textbf{Meaning-based search.} Older search compared exact words. Dense
passage retrieval~\cite{karpukhin2020dpr} instead compares meaning, which
finds the right text even when the words differ. Sentence-BERT
~\cite{reimers2019sbert} made it quick to turn short text into numbers for
this kind of comparison. A second model, a cross-encoder, can re-check
the shortlisted text against the question to keep only the best matches.

\textbf{Using AI models to write answers.} Current AI models read given
text and answer from it quite well. But unless they are clearly told to
use only that text, they tend to add their own knowledge. For exam study
this is a problem, because only the set book should be used.

VernaSolver stands apart in three ways: the page number is always shown
to the student, the system refuses when the book lacks the answer, and it
is built for one-book study in more than one language on ordinary
hardware.

% =====================================================================
%  III. EXISTING VS PROPOSED SYSTEM
% =====================================================================
\section{Existing and Proposed System}
Table~\ref{tab:compare} compares a normal AI chat tool with VernaSolver.
A normal tool is general and easy to use, but it cannot promise that an
answer is from the student's book, and it often cannot be checked.
VernaSolver trades some generality for trust: it only works on books that
have been added, but every answer it gives can be traced to a page in
those books.

\begin{table}[!t]
\caption{Normal Chat Tool vs. VernaSolver}
\label{tab:compare}
\centering
\renewcommand{\arraystretch}{1.25}
\begin{tabular}{@{}p{2.5cm}p{2.4cm}p{2.4cm}@{}}
\toprule
\textbf{Point} & \textbf{Normal Chat Tool} & \textbf{VernaSolver} \\
\midrule
Source of answer & Training data & The student's book \\
Page number      & Often wrong/missing & Shown and real \\
Made-up answers  & Common & Avoided; says ``not found'' \\
Syllabus match   & None & Limited to chosen book \\
Languages        & Mostly English & English, Hindi, Marathi \\
Study aids       & General & Quiz/flashcards from book \\
\bottomrule
\end{tabular}
\end{table}

% =====================================================================
%  IV. SYSTEM DESIGN
% =====================================================================
\section{System Design}
VernaSolver runs in three steps: reading the book, searching, and
answering. Fig.~\ref{fig:flow} shows the flow of data.

% ---- Flow figure (native TikZ) ----
\begin{figure}[!t]
\centering
\resizebox{0.92\columnwidth}{!}{%
\begin{tikzpicture}[
  node distance=6mm,
  font=\footnotesize,
  step/.style={rectangle, rounded corners=2pt, draw=black, very thin,
               minimum height=8mm, minimum width=6.4cm, align=center,
               fill=blue!6},
  io/.style={rectangle, rounded corners=2pt, draw=black, very thin,
             minimum height=8mm, minimum width=6.4cm, align=center,
             fill=green!10},
  arr/.style={-{Stealth[length=2mm]}, thin}
]
\node[io]                       (a) {\textbf{Input:} Student uploads a PDF book};
\node[step, below=of a]         (b) {\textbf{Step 1 -- Read:} split into parts, turn into numbers, save};
\node[step, below=of b]         (c) {\textbf{Step 2 -- Search:} find best parts, re-order with 2nd model};
\node[step, below=of c]         (d) {\textbf{Step 3 -- Answer:} AI writes reply using only those parts};
\node[io, below=of d]           (e) {\textbf{Output:} Answer with page numbers};
\draw[arr] (a) -- (b);
\draw[arr] (b) -- (c);
\draw[arr] (c) -- (d);
\draw[arr] (d) -- (e);
\end{tikzpicture}%
}
\caption{The three main steps of VernaSolver, from uploading a book to
getting an answer with page numbers.}
\label{fig:flow}
\end{figure}

\subsection{Reading the Book}
When a PDF is uploaded, the system reads it page by page with the PyMuPDF
library. Reading this way keeps a record of which text is on which page,
so the right page can be shown later. The text is split into parts of
about 400 words, with a 40-word overlap so an idea that crosses a boundary
is not cut off. Each part is turned into a list of 384 numbers using a
small model named \texttt{all-MiniLM-L6-v2}, which runs on a normal
computer. These numbers, with the book title, author, subject, and page,
are stored in a ChromaDB database.

\subsection{Searching}
The question is turned into numbers with the same model and matched
against the stored parts to find the closest ones. A second model
(\texttt{ms-marco-MiniLM-L-6-v2}) then looks at the question and each
shortlisted part together and keeps the best ones. If a student asks a
short follow-up like ``and its uses?'', the system uses the earlier chat
to rewrite it into a full question before searching.

\subsection{Answering}
The chosen parts and their page numbers go to an AI model (Anthropic
Claude, with OpenAI GPT-4o-mini as a backup). The model is told to use
only these parts, to copy maths steps as written in the book instead of
solving them itself, and to clearly say when the answer is not present.
The reply appears word by word on screen using Server-Sent Events (SSE),
and the page numbers are listed below it.

% =====================================================================
%  V. KEEPING ANSWERS TIED TO THE BOOK
% =====================================================================
\section{Keeping Answers Tied to the Book}
Three steps keep the answers trustworthy. First, the instruction to the
AI model says to use only the parts taken from the book and to say so
when the book lacks the answer; this makes the page number a real promise.
Second, because each part keeps its page number through the whole process,
every answer can be traced to a page, and that page can even be shown from
the PDF with the used text marked. Third, the search runs on the local
computer while the AI model only writes the final answer; this keeps cost
and waiting time low and makes sure the model reads from the book instead
of from memory.

For quizzes and flashcards, the same search step gathers related text and
the model is asked to reply in a fixed JSON form. To stop it from writing
a normal paragraph, its reply is started with an opening brace so it has
to continue as JSON.

% =====================================================================
%  VI. IMPLEMENTATION
% =====================================================================
\section{Implementation}
VernaSolver is written in Python. Table~\ref{tab:stack} lists the main
tools. The server is built with FastAPI and offers web addresses for
chat, book management, uploads, quizzes, flashcards, and login. A
command-line tool built with Click does the same core work and is useful
for testing. The web page uses plain HTML, CSS, and JavaScript and shows
the answer as it streams in. Accounts are kept in SQLite, with passwords
stored in a hashed form (PBKDF2-SHA256) and login handled with cookies.

\begin{table}[!t]
\caption{Tools Used in VernaSolver}
\label{tab:stack}
\centering
\renewcommand{\arraystretch}{1.25}
\begin{tabular}{@{}ll@{}}
\toprule
\textbf{Part} & \textbf{Tool} \\
\midrule
Server framework & FastAPI (Python) \\
Reading PDF      & PyMuPDF (fitz) \\
Text-to-numbers  & all-MiniLM-L6-v2 (local) \\
Re-checking      & ms-marco-MiniLM-L-6-v2 \\
Search database  & ChromaDB \\
Main AI model    & Anthropic Claude \\
Backup AI model  & OpenAI GPT-4o-mini \\
Live replies     & Server-Sent Events (SSE) \\
Login \& accounts & SQLite + PBKDF2-SHA256 \\
Web page         & HTML / CSS / JavaScript \\
Command-line tool & Click \\
\bottomrule
\end{tabular}
\end{table}

\subsection{Main Settings}
Each part of the book is about 400 words, with a 40-word overlap and a
30-word minimum. The search takes 12 close matches and keeps the best 5
after re-checking. Up to 6 recent question-and-answer turns are kept for
context. These values give good search results without sending too much
text to the AI model at once.

% =====================================================================
%  VII. RESULTS & DISCUSSION
% =====================================================================
\section{Results and Discussion}
We tested a working version with real course books, including
\textit{Software Engineering} by Pressman and \textit{Concepts of
Physics} by H.C. Verma. The system answered concept and definition
questions correctly and showed the right page numbers. When a book did
not have the answer, it said so instead of guessing. It also produced
quizzes and flashcards on chosen topics, and replies appeared smoothly as
they were written. We noticed two things that helped the most: telling the
model to stay inside the book kept answers faithful to the source, and the
re-checking step removed many weak matches before the answer was written.
Rewriting short follow-up questions also made longer chats work much
better.

% =====================================================================
%  VIII. CONCLUSION & FUTURE SCOPE
% =====================================================================
\section{Conclusion and Future Scope}
This paper presented VernaSolver, a study tool that answers only from a
student's chosen book and shows a page number for every answer. By keeping
the search on a normal computer, using the AI model only to write the
reply, and telling it to stay inside the book, the system avoids the
made-up answers and the lack of trust seen in normal chat tools. It also
supports Indian languages and makes quizzes and flashcards.

In future work we plan to add a step that reads text from scanned books,
use a search model that handles Hindi and Marathi better, run proper tests
to measure answer and page-number accuracy, and add per-user limits and
saved chat history so the tool is ready for wider use.

% =====================================================================
%  REFERENCES
% =====================================================================
\begin{thebibliography}{00}

\bibitem{lewis2020rag}
P. Lewis \textit{et al.}, ``Retrieval-augmented generation for
knowledge-intensive NLP tasks,'' in \textit{Advances in Neural
Information Processing Systems (NeurIPS)}, 2020.

\bibitem{karpukhin2020dpr}
V. Karpukhin \textit{et al.}, ``Dense passage retrieval for open-domain
question answering,'' in \textit{Proc. EMNLP}, 2020, pp.~6769--6781.

\bibitem{reimers2019sbert}
N. Reimers and I. Gurevych, ``Sentence-BERT: Sentence embeddings using
Siamese BERT-networks,'' in \textit{Proc. EMNLP-IJCNLP}, 2019,
pp.~3982--3992.

\bibitem{vaswani2017attention}
A. Vaswani \textit{et al.}, ``Attention is all you need,'' in
\textit{Advances in Neural Information Processing Systems (NeurIPS)},
2017.

\bibitem{chromadb}
Chroma, ``ChromaDB: the open-source embedding database,''
\url{https://www.trychroma.com}. (Accessed 2026).

\end{thebibliography}

\end{document}
